% GENERATED FILE. Edit canonical lessons, then run npm run generate:books.
% canonical-source-hash: fnv1a64:dbf54679853ddccf
% canonical-lessons: FA-C497-ketf, FA-C497-mafsal, FA-C497-naf, FA-C497-sotun-e-faqarat, FA-C497-jomjome

\chapter{Bones and Joints}
\label{ch:fa-bones-and-joints}
\hlchaptermodality{497}

\begin{chapteropening}
\textbf{By the end of this chapter:} \emph{I can talk about bones and joints.}

Complete the last lesson of chapter 497: I can talk about bones and joints.
\end{chapteropening}

\section[\texorpdfstring{\fa{کتف} (ketf)}{کتف (ketf)}]{\fa{کتف} (ketf) — a shoulder blade; shoulder}

\label{lesson:FA-C497-ketf}



\begin{quote}
Before the new one: say the Persian for an eyelid, then the Persian for the gum.
\end{quote}

\subsection*{You'll want to know: \fa{کتف}}
\textbf{\fa{کتف}} (\emph{ketf}) — \textquotedblleft{}a shoulder blade; shoulder\textquotedblright{}.

\begin{quote}
\fa{او} \fa{کیف} \fa{سنگینی} \fa{روی} \fa{کتفش} \fa{داشت}.
\end{quote}

\begin{quote}
\emph{u kif-e sangini ru-ye ketfash dâsht.}
\end{quote}

\begin{quote}
She had a heavy bag on her shoulder.
\end{quote}

\emph{shâne} is the shoulder and \emph{ketf} the shoulder blade, but in speech they overlap.

\begin{cousinweb}
\emph{ketf} comes from Arabic \emph{kitf}, 'shoulder' (EtymDB).
\end{cousinweb}

\begin{grammarlens}[title={What you've built}]
One more word for telling, explaining and describing at B1.
\end{grammarlens}

\subsection*{Your turn}
\begin{itemize}
\raggedright
  \item \emph{Say it:} \emph{ketf}
  \item \emph{Say it:} \emph{ketf}, once more
  \item \emph{Say it:} a sentence of your own with \emph{ketf}
  \item \emph{Recall:} say the Persian for an eyelid, then the Persian for the gum, then say \emph{ketf} again
\end{itemize}

\begin{tcolorbox}[breakable,colback=teal!4,colframe=teal!35!black,title={Before you move on}]
What does \fa{کتف} mean? (\textquotedblleft{}A shoulder blade; shoulder\textquotedblright{}.) Say it once more.
\end{tcolorbox}
\section[\texorpdfstring{\fa{مفصل} (mafsal)}{مفصل (mafsal)}]{\fa{مفصل} (mafsal) — a joint (of the body)}

\label{lesson:FA-C497-mafsal}



\begin{quote}
Before the new one: say the Persian for the gum, then the Persian for a shoulder blade.
\end{quote}

\subsection*{You'll want to know: \fa{مفصل}}
\textbf{\fa{مفصل}} (\emph{mafsal}) — \textquotedblleft{}a joint (of the body)\textquotedblright{}.

\begin{quote}
\fa{مادربزرگ} \fa{در} \fa{زمستان} \fa{درد} \fa{مفصل} \fa{دارد}.
\end{quote}

\begin{quote}
\emph{mâdarbozorg dar zemestân dard-e mafsal dârad.}
\end{quote}

\begin{quote}
Grandmother has joint pain in winter.
\end{quote}

Plural \emph{mafâsel} (Arabic). The same letters read \emph{mofassal} mean 'detailed'.

\begin{cousinweb}
\emph{mafsal} is a loan, usually traced to Arabic \emph{mafṣil}, 'joint'. EtymDB's entry under this spelling is the other reading, \emph{mufaṣṣal}, 'divided' (EtymDB), which Persian keeps as \emph{mofassal}, 'detailed'.
\end{cousinweb}

\begin{grammarlens}[title={What you've built}]
One more word for telling, explaining and describing at B1.
\end{grammarlens}

\subsection*{Your turn}
\begin{itemize}
\raggedright
  \item \emph{Say it:} \emph{mafsal}
  \item \emph{Say it:} \emph{mafsal}, once more
  \item \emph{Say it:} a sentence of your own with \emph{mafsal}
  \item \emph{Recall:} say the Persian for the gum, then the Persian for a shoulder blade, then say \emph{mafsal} again
\end{itemize}

\begin{tcolorbox}[breakable,colback=teal!4,colframe=teal!35!black,title={Before you move on}]
What does \fa{مفصل} mean? (\textquotedblleft{}A joint (of the body)\textquotedblright{}.) Say it once more.
\end{tcolorbox}
\section[\texorpdfstring{\fa{ناف} (nâf)}{ناف (nâf)}]{\fa{ناف} (nâf) — the navel}

\label{lesson:FA-C497-naf}



\begin{quote}
Before the new one: say the Persian for a shoulder blade, then the Persian for a joint.
\end{quote}

\subsection*{You'll want to know: \fa{ناف}}
\textbf{\fa{ناف}} (\emph{nâf}) — \textquotedblleft{}the navel\textquotedblright{}.

\begin{quote}
\fa{پرستار} \fa{بند} \fa{ناف} \fa{بچه} \fa{را} \fa{بست}.
\end{quote}

\begin{quote}
\emph{parastâr band-e nâf-e bachche râ bast.}
\end{quote}

\begin{quote}
The nurse tied the baby's umbilical cord.
\end{quote}

\emph{band-e nâf} — the umbilical cord; \emph{nâf-e shahr} — the very heart of a town.

\begin{cousinweb}
\emph{nâf} is inherited: EtymDB traces it to a Proto-Indo-European word for 'navel', and lists Sanskrit \emph{nābhi} and Avestan as cousins (EtymDB).
\end{cousinweb}

\begin{grammarlens}[title={What you've built}]
One more word for telling, explaining and describing at B1.
\end{grammarlens}

\subsection*{Your turn}
\begin{itemize}
\raggedright
  \item \emph{Say it:} \emph{nâf}
  \item \emph{Say it:} \emph{nâf}, once more
  \item \emph{Say it:} a sentence of your own with \emph{nâf}
  \item \emph{Recall:} say the Persian for a shoulder blade, then the Persian for a joint, then say \emph{nâf} again
\end{itemize}

\begin{tcolorbox}[breakable,colback=teal!4,colframe=teal!35!black,title={Before you move on}]
What does \fa{ناف} mean? (\textquotedblleft{}The navel\textquotedblright{}.) Say it once more.
\end{tcolorbox}
\section[\texorpdfstring{\fa{ستون} \fa{فقرات} (sotun-e faqârât)}{ستون فقرات (sotun-e faqârât)}]{\fa{ستون} \fa{فقرات} (sotun-e faqârât) — the spine, backbone}

\label{lesson:FA-C497-sotun-e-faqarat}



\begin{quote}
Before the new one: say the Persian for a joint, then the Persian for the navel.
\end{quote}

\subsection*{You'll want to know: \fa{ستون} \fa{فقرات}}
\textbf{\fa{ستون} \fa{فقرات}} (\emph{sotun-e faqârât}) — \textquotedblleft{}the spine, backbone\textquotedblright{}.

\begin{quote}
\fa{تصادف} \fa{به} \fa{ستون} \fa{فقرات} \fa{او} \fa{آسیب} \fa{زد}.
\end{quote}

\begin{quote}
\emph{tasâdof be sotun-e faqârât-e u âsib zad.}
\end{quote}

\begin{quote}
The accident injured his spine.
\end{quote}

\emph{sotun} 'pillar' + \emph{faqârât} 'vertebrae' (an Arabic plural): the pillar of vertebrae.

\begin{cousinweb}
\emph{sotun-e faqârât} is built in Persian on inherited \emph{sotun}, 'pillar', usually traced to Middle Persian \emph{stūn}, and \emph{faqârât}, 'vertebrae', a loan usually traced to the Arabic plural \emph{faqārāt}.
\end{cousinweb}

\begin{grammarlens}[title={What you've built}]
One more word for telling, explaining and describing at B1.
\end{grammarlens}

\subsection*{Your turn}
\begin{itemize}
\raggedright
  \item \emph{Say it:} \emph{sotun-e faqârât}
  \item \emph{Say it:} \emph{sotun-e faqârât}, once more
  \item \emph{Say it:} a sentence of your own with \emph{sotun-e faqârât}
  \item \emph{Recall:} say the Persian for a joint, then the Persian for the navel, then say \emph{sotun-e faqârât} again
\end{itemize}

\begin{tcolorbox}[breakable,colback=teal!4,colframe=teal!35!black,title={Before you move on}]
What does \fa{ستون} \fa{فقرات} mean? (\textquotedblleft{}The spine, backbone\textquotedblright{}.) Say it once more.
\end{tcolorbox}
\section[\texorpdfstring{\fa{جمجمه} (jomjome)}{جمجمه (jomjome)}]{\fa{جمجمه} (jomjome) — a skull}

\label{lesson:FA-C497-jomjome}



\begin{quote}
Before the new one: say the Persian for the navel, then the Persian for the spine.
\end{quote}

\subsection*{You'll want to know: \fa{جمجمه}}
\textbf{\fa{جمجمه}} (\emph{jomjome}) — \textquotedblleft{}a skull\textquotedblright{}.

\begin{quote}
\fa{در} \fa{موزه} \fa{یک} \fa{جمجمه} \fa{خیلی} \fa{قدیمی} \fa{بود}.
\end{quote}

\begin{quote}
\emph{dar muze yek jomjome-ye kheyli qadimi bud.}
\end{quote}

\begin{quote}
In the museum there was a very old skull.
\end{quote}

Plural \emph{jomjomehâ}.

\begin{cousinweb}
\emph{jomjome} was borrowed from Arabic \emph{jumjuma}, 'skull' (EtymDB).
\end{cousinweb}

\begin{grammarlens}[title={What you've built}]
One more word for telling, explaining and describing at B1.
\end{grammarlens}

\subsection*{Your turn}
\begin{itemize}
\raggedright
  \item \emph{Say it:} \emph{jomjome}
  \item \emph{Say it:} \emph{jomjome}, once more
  \item \emph{Say it:} a sentence of your own with \emph{jomjome}
  \item \emph{Recall:} say the Persian for the navel, then the Persian for the spine, then say \emph{jomjome} again
\end{itemize}

\begin{tcolorbox}[breakable,colback=teal!4,colframe=teal!35!black,title={Before you move on}]
What does \fa{جمجمه} mean? (\textquotedblleft{}A skull\textquotedblright{}.) Say it once more.
\end{tcolorbox}
